\documentclass[UTF8,11pt,a4paper,fontset=fandol]{ctexart}
\usepackage[margin=20mm]{geometry}
\usepackage{amsmath,amssymb,booktabs,array,tabularx,enumitem,setspace,xcolor,hyperref,xurl,fancyhdr}
\setstretch{1.08}
\setlength{\emergencystretch}{3em}
\setlength{\headheight}{15pt}
\setlength{\tabcolsep}{5pt}
\renewcommand{\arraystretch}{1.2}
\setlist{nosep,leftmargin=2em}
\hypersetup{colorlinks=true,linkcolor=blue!45!black,urlcolor=blue!45!black}
\providecommand{\WorkloadName}{工作负载}
\pagestyle{fancy}\fancyhf{}
\fancyhead[L]{\small Table II(b)：算子级 RI}
\fancyhead[R]{\small\WorkloadName}
\fancyfoot[C]{\thepage}
\newcommand{\qs}{Q_{\mathrm S}}
\newcommand{\qr}{Q_{\mathrm R}}
\newcommand{\RI}{\mathrm{RI}}
\newcommand{\para}[1]{\par\smallskip\noindent\textbf{#1}\quad}
\newcommand{\ModelTableSetup}{\renewcommand{\arraystretch}{1.2}\setlength{\tabcolsep}{4pt}}
\newcommand{\ResultTableSetup}{\renewcommand{\arraystretch}{1.45}\setlength{\tabcolsep}{4pt}}
\newcommand{\DisplayNote}{\par\smallskip{\noindent\footnotesize $1\mathrm{K}=1024$，$1\mathrm{M}=1024^2$。RI 至少为 1 时保留一位小数，小于 1 时保留三位有效数字；精确值见机器数据。\par}\smallskip}
\author{}\date{2026 年 10 月 2 日}

% Compact, identical title block keeps the mathematical introduction on one page.
\makeatletter
\renewcommand{\maketitle}{\thispagestyle{plain}\begin{center}
{\LARGE\bfseries\@title\par}\vspace{8pt}{\normalsize\@date\par}
\end{center}\vspace{3pt}}
\makeatother
